\part{Stand}

\chapter{Wo die Vorhersage endet}

\section{Der Leitfall, ein letztes Mal}

\begin{merksatz}
\boxtitel{Was der Band aus einer Herzzelle gemacht hat}
\begin{center}
{\small
\renewcommand{\arraystretch}{1.46}
\begin{tabular}{@{}P{1.2cm}P{3.4cm}P{3.8cm}@{}}
\toprule
\textbf{Kap.} & \textbf{Die Zelle} & \textbf{Der Satz}\\
\midrule
$1$--$2$ & Ruhe bei $(-1{,}20;\ -0{,}62)$, Schwelle $0{,}555$ & Fluss, Fixpunkt, Spur und Determinante\\
$3$ & Schlag alle $39{,}5$ bei $I = 0{,}5$ & Grenzzyklus, Rückkehrabbildung\\
$4$--$5$ & Takt entsteht bei $I = 0{,}3313$, subkritisch & Sattel-Knoten, Hopf, Hysterese $0{,}3245$--$0{,}3313$\\
$6$ & $1\!:\!1$, $2\!:\!1$, $3\!:\!2$ unter Reizen & Kreisabbildung, Arnold-Zungen\\
$7$ & Alternans bei $\mathrm{BCL} = 326$ ms & Flip-Verzweigung, Feigenbaum\\
$9$ & $T_p = 18$: $\lambda \approx 0$ --- nicht chaotisch & Ljapunow-Exponent\\
$10$--$12$ & aus einem Kanal; Test gegen Zufall & Takens, Dimension, Surrogate\\
$13$ & Welle $c = 1{,}03$, Spirale $T = 38{,}7$ & Front, Eikonal, Wellenlänge\\
$14$ & tausend Zellen: Dichte statt Zahl & Kuramoto, extensives Chaos\\
\bottomrule
\end{tabular}}
\end{center}

\begin{center}
\begin{minipage}{0.9\textwidth}
\textbf{Von zwei Gleichungen bis zum Gewebe hat jede Zahl des Bandes an derselben Zelle gestanden. Was sich geändert hat, ist nicht die Zelle, sondern die Frage: erst \enquote{was tut sie}, dann \enquote{wann ändert sie sich}, dann \enquote{kann man es ihr ansehen}, zuletzt \enquote{was tun Milliarden davon}.}
\end{minipage}
\end{center}
\end{merksatz}

\section{Vier Muster}

\begin{merksatz}
\boxtitel{Muster 1 --- die Grenze ist ein Verhältnis}
\begin{center}
{\small
\renewcommand{\arraystretch}{1.5}
\begin{tabular}{@{}P{1.4cm}P{7.0cm}@{}}
\toprule
Kap.~8 & $10^{-8}$ wird nach $t = 20$ zu $4$.\\
Kap.~9 & $t_{\max} = \ln(\Delta/\delta_0)/\lambda_1$; eine Stelle mehr bringt $\ln10/\lambda_1$.\\
Kap.~9 & Wetter: zwei Wochen, Sonnensystem: zehn Millionen Jahre --- dieselbe Formel.\\
\bottomrule
\end{tabular}}
\end{center}

\begin{center}
\begin{minipage}{0.88\textwidth}
\kursiv{\textbf{Der Titel des Bandes hat eine Formel, und sie hat zwei Eingänge}: die Genauigkeit des Anfangs und die Rate des Vergessens. Die erste kann man verbessern --- logarithmisch wirksam; die zweite ist eine Eigenschaft der Gleichungen. Wer die Vorhersage verlängern will, muss nicht besser messen, sondern das System wechseln: Ensembles statt Verläufe (Kapitel 9), Statistik des Attraktors statt Bahn (Kapitel 11), oder --- die Antwort der Medizin --- nicht vorhersagen, sondern verhindern: die Wellenlänge verlängern, bis kein Kreis mehr passt.}
\end{minipage}
\end{center}
\end{merksatz}

\begin{merksatz}
\boxtitel{Muster 2 --- unregelmäßig ist nicht chaotisch}
\begin{center}
{\small
\renewcommand{\arraystretch}{1.5}
\begin{tabular}{@{}P{1.4cm}P{7.0cm}@{}}
\toprule
Kap.~6, 9 & Die Zone bei $T_p = 18$: sieht chaotisch aus, $\lambda \approx 0$, quasiperiodisch.\\
Kap.~8 & Der Sattel trennt Verläufe, ist aber kein Chaos: keine Rückkehr.\\
Kap.~11 & Farbiges Rauschen hat eine Dimension.\\
Kap.~13 & Die Spirale ist periodisch; das EKG der Tachykardie regelmäßig.\\
\bottomrule
\end{tabular}}
\end{center}

\begin{center}
\begin{minipage}{0.88\textwidth}
\kursiv{\textbf{Chaos ist ein Fachbegriff mit drei Bedingungen --- beschränkt, empfindlich, wiederkehrend ---, und das Auge prüft keine davon.} Der Band hat an vier Stellen gezeigt, wie Unregelmäßigkeit ohne Chaos entsteht: Quasiperiodizität, Instabilität, Zufall mit Gedächtnis, Ordnung auf einer Skala, die man nicht sieht. Umgekehrt hat die Alternans-Kaskade gezeigt, wie Chaos aus vollkommener Regelmäßigkeit entsteht, in Schritten, die man zählen kann. Die Frage \enquote{ist das Chaos?} ist deshalb nie eine Frage des Eindrucks, sondern eine der Messung --- und die Messung braucht den Test aus Kapitel 12.}
\end{minipage}
\end{center}
\end{merksatz}

\begin{merksatz}
\boxtitel{Muster 3 --- das Modell ist nicht der Patient}
\begin{center}
{\small
\renewcommand{\arraystretch}{1.5}
\begin{tabular}{@{}P{1.4cm}P{7.0cm}@{}}
\toprule
Kap.~5 & Die Hopf-Verzweigung ist subkritisch --- im Modell; das Lehrbuch sagt superkritisch.\\
Kap.~9 & FitzHugh--Nagumo mit Stößen wird nicht chaotisch; die Zellen von Guevara und Glass wurden es.\\
Kap.~13 & Die Spirale bricht nicht --- flache Restitution; Herzen brechen.\\
Kap.~13 & Die Frontformel gibt $2{,}0$, die Simulation $1{,}03$: $\varepsilon$ ist nicht klein.\\
\bottomrule
\end{tabular}}
\end{center}

\begin{center}
\begin{minipage}{0.88\textwidth}
\kursiv{\textbf{Viermal hat das Modell etwas anderes getan als erwartet, und jedes Mal stand es im Text.} Das ist keine Schwäche von FitzHugh--Nagumo, sondern sein Zweck: Zwei Gleichungen, die alles zeigen, was eine erregbare Zelle kann --- und an denen man sieht, wo die Zelle mehr ist. Die Modelle, mit denen 2026 gerechnet wird (\Person{ten Tusscher} 2004, \Person{O'Hara}--\Person{Rudy} 2011), haben vierzig Veränderliche, und ihre Restitution bricht Spiralen; aber die Begriffe, mit denen man ihre Ergebnisse liest --- Schwelle, Zyklus, Zunge, Flip, Wellenlänge, Kern ---, sind die der zwei Gleichungen.}
\end{minipage}
\end{center}
\end{merksatz}

\begin{merksatz}
\boxtitel{Muster 4 --- Dehnen und Falten}
\begin{center}
{\small
\renewcommand{\arraystretch}{1.5}
\begin{tabular}{@{}P{1.4cm}P{7.0cm}@{}}
\toprule
Kap.~7 & Der Buckel der logistischen Abbildung: Dehnen (Steigung $> 1$), Falten (das Maximum).\\
Kap.~8 & Lorenz: Dehnen $\mathrm e^{0{,}906}$, Zusammenziehen $\mathrm e^{-14{,}57}$, Falten beim Flügelwechsel.\\
Kap.~11 & Die $0{,}06$ in $2{,}06$: die Cantor-Menge quer zum Blatt.\\
Kap.~13 & Der Wellenbruch: Dehnen der Spirale durch Alternans, Falten durch die Refraktärzone.\\
\bottomrule
\end{tabular}}
\end{center}

\begin{center}
\begin{minipage}{0.88\textwidth}
\kursiv{\textbf{Ein einziger Mechanismus steht hinter jedem Chaos des Bandes}: Nachbarn auseinanderziehen, dann den Raum zurückfalten, damit es beschränkt bleibt. Smales Hufeisen von 1967 ist seine Reinform (Auslassung 17), und alles --- Exponent, Dimension, Entropie --- sind Zahlen, die sagen, wie stark gedehnt und wie fein gefaltet wird. Wer in einem neuen System Chaos vermutet, suche zuerst, wo es dehnt und wo es faltet; findet er nur eines, ist es keines.}
\end{minipage}
\end{center}
\end{merksatz}

\section{Was offen ist}

\begin{merksatz}
\boxtitel{Stand bei Redaktionsschluss, August 2026}
\begin{center}
{\small
\renewcommand{\arraystretch}{1.5}
\begin{tabular}{@{}P{3.2cm}P{5.4cm}@{}}
\toprule
\textbf{Gelernte Wettervorhersage} & Seit 2023 (GraphCast, Pangu) schlagen neuronale Netze, die aus vierzig Jahren Reanalysen gelernt haben, die physikalischen Modelle bei zehn Tagen; seit 2025 läuft am EZMW ein solches Modell im Betrieb, Ensembles aus Diffusionsmodellen (GenCast 2024) geben Verteilungen. Die zwei Wochen aus Kapitel 9 haben sie nicht verschoben --- sie erreichen sie billiger.\\
\textbf{Reservoir Computing} & Ein zufälliges Netz als Einbettung sagt Lorenz über acht Ljapunow-Zeiten voraus (Pathak 2018) und reproduziert das Spektrum räumlich ausgedehnter Systeme; warum es so gut geht, ist nur teilweise verstanden.\\
\textbf{Der digitale Zwilling} & Patientenspezifische Herzmodelle aus MRT und Elektrophysiologie (Trayanova 2018, Niederer 2019) sagen den Ort der Kammertachykardie voraus und leiten die Ablation; erste Studien, kein Standard.\\
\textbf{Rotoren beim Patienten} & Ob Vorhofflimmern von wenigen stabilen Spiralen getrieben wird (Narayan 2012) oder von vielen zerfallenden Wellen (Allessie), ist nach negativen Studien zur Rotorablation (REAFFIRM 2019) offen --- und mit Kapitel 14 vielleicht falsch gestellt.\\
\textbf{Ist das Herz chaotisch?} & Nichtlinear ja, chaotisch unentscheidbar (Glass 2009); die Frage hat die Messgrößen hinterlassen, nicht die Antwort.\\
\bottomrule
\end{tabular}}
\end{center}

\begin{center}
\begin{minipage}{0.88\textwidth}
\kursiv{\textbf{Das Bemerkenswerteste an den gelernten Wettermodellen ist, was sie nicht geändert haben.} Sie kennen keine Gleichung, nur Daten; sie rechnen in Minuten statt Stunden; sie schlagen das beste physikalische Modell in fast jeder Größe --- und ihr Vorhersagehorizont liegt dort, wo \Person{Lorenz} ihn 1969 gesetzt hat. Das ist Satz 9.4 als Experiment: Die Grenze liegt in den Gleichungen der Atmosphäre, nicht in unserer Fähigkeit, sie zu lösen, und ein Modell, das die Gleichungen aus Daten lernt, erbt die Grenze mit. Was sich verschiebt, ist die Form der Vorhersage: von einem Verlauf zu einer Verteilung, von \enquote{es regnet am Dienstag} zu \enquote{in $70$ Prozent der Welten, die mit unserem Wissen verträglich sind}. Band 24 hat gesagt, was eine solche Aussage ist; dieser Band hat gesagt, warum es keine bessere gibt.}
\end{minipage}
\end{center}
\end{merksatz}

\section{Sechs Sätze, die bleiben}

\begin{merksatz}
\boxtitel{Die Bilanz}
\begin{center}
{\small
\renewcommand{\arraystretch}{1.5}
\begin{tabular}{@{}P{2.0cm}P{6.6cm}@{}}
\toprule
\textbf{Satz 2.3} & Spur und Determinante entscheiden über den Fixpunkt --- zwei Zahlen für jede Ruhe.\\
\textbf{Satz 5.1} & Hopf: Der Takt entsteht aus der Ruhe, und sein Vorzeichen entscheidet, ob sanft oder mit Sprung.\\
\textbf{Satz 7.3} & Feigenbaum: $4{,}669$ für jeden Buckel, gemessen in Helium und Dioden.\\
\textbf{Satz 9.4} & Der Horizont ist $\ln(\Delta/\delta_0)/\lambda_1$ --- logarithmisch in der Genauigkeit.\\
\textbf{Satz 10.2} & Takens: Ein Kanal genügt, weil die Vergangenheit die übrigen Veränderlichen enthält.\\
\textbf{Satz 13.3} & Die Wellenlänge $c\,T_{\mathrm{ref}}$ entscheidet, ob ein Kreis passt --- und damit über Reentry.\\
\bottomrule
\end{tabular}}
\end{center}

\begin{center}
\begin{minipage}{0.9\textwidth}
\textbf{Dazu ein Test (Satz 12.2), der sagt, wann die Sätze überhaupt anwendbar sind. Ohne ihn sind die anderen sechs Werkzeuge, die an Rauschen genauso gut funktionieren.}
\end{minipage}
\end{center}
\end{merksatz}

\begin{merksatz}
\boxtitel{Ein Schlusswort}
\begin{center}
\begin{minipage}{0.9\textwidth}
\kursiv{Band 27 hat gezeigt, was der Rechner vom Zufall weiß: dass er ihn nachbauen kann, mit Fehlern, die wie $1/\sqrt n$ fallen. Dieser Band hat gezeigt, dass der Rechner auch ohne Zufall an eine Grenze kommt --- dass drei glatte Gleichungen eine Bahn erzeugen, deren achte Stelle nach zwanzig Umläufen die erste ist, und dass keine Genauigkeit der Welt das ändert, sondern nur verschiebt, um einen festen Betrag je Stelle. \textbf{Das ist keine Niederlage der Vorhersage, sondern ihre Vermessung.} Man weiß, wo sie endet, und man weiß, warum; man kann die Rate messen, an der sie endet, aus einem einzigen Kanal; man kann prüfen, ob sie an Chaos endet oder an Zufall; und man kann, wo sie endet, etwas anderes vorhersagen --- die Verteilung, den Attraktor, die Dichte. Die Herzzelle, die den Band begleitet hat, ist der Beweis, dass das kein Trost ist, sondern Medizin: Die Wellenlänge, die über Flimmern entscheidet, ist eine Vorhersage, die hält. Wo die Gleichungen selbst unendlich viele Veränderliche haben --- in der Turbulenz ---, beginnt die Physik, und dorthin führt eine andere Reihe.}
\end{minipage}
\end{center}
\end{merksatz}
